\chapter{Introduction}\label{chap:introducao}
% Funnel structure — start broad, end at the specific research gap.
%
% 1. Technological/academic/scientific context
%    - 5G high data rates; 6G even higher demands
%    - MIMO and Massive MIMO as key enabling technologies
%    - High computational cost of Massive MIMO signal processing
%
% 2. Narrowing toward the core problem
%    - Linear equalizers (LMMSE) achieve near-ML spectral efficiency
%    - Closed-form solution requires matrix inversion: \mathcal{O}(N^3) complexity
%    - Under colored noise (e.g., antenna coupling), \mathbf{R}_n^{-1} cannot be
%      avoided or trivially decomposed as in AWGN
%    - Real-time detection requirement creates a processing bottleneck
%
% 3. Research gap and objective
%    - Decentralized architectures reduce per-cluster matrix size, but
%      inversion cost remains relevant
%    - Zhao et al. (2025) propose BCD-MMSE, the first decentralized LMMSE
%      equalizer under colored noise; its per-DU coherence-block inversion
%      (Eq. 30) has no hardware implementation
%    - This work proposes an FPGA module for that inversion step,
%      evaluated on BER, LUT usage, and power consumption

The evolution toward fifth-generation (5G) and sixth-generation (6G) wireless
systems demands ever-increasing spectral efficiency and data throughput, driving
the adoption of Massive Multi-User Multiple-Input Multiple-Output (MU-MIMO) as
a key enabling technology~\cite{decentralized2017}. By equipping base stations
with hundreds of antenna elements serving tens of users simultaneously, Massive
MU-MIMO achieves fine-grained spatial multiplexing that small-scale MIMO systems
cannot provide. Among the available uplink detection strategies, Linear Minimum
Mean-Square Error (LMMSE) equalization is widely adopted for its near-optimal
performance~\cite{zhao2025jsac}. Its closed-form solution, however, requires a
matrix inversion whose complexity grows cubically with the number of antennas,
and for a base station with 256 antennas operating at 80~MHz bandwidth with
12-bit converters, the raw fronthaul data rate approaches 1~Tbps, far beyond
the capacity of existing interconnect standards such as the Common Public Radio
Interface (CPRI)~\cite{li2019cd}.

Decentralized Baseband Processing (DBP) architectures address both bottlenecks
by partitioning the base-station antenna array into clusters of $M_c$ antennas, each served by a local
Distributed Unit (DU) that performs baseband processing independently. This
reduces per-cluster matrix dimensions and eliminates the need to transfer raw
antenna data to a centralized processor~\cite{jeon2019}. Existing decentralized
LMMSE methods, however, almost universally assume that receiver noise follows an ideal
Additive White Gaussian Noise (AWGN) model with a diagonal covariance matrix,
an assumption that typically does not hold in practice, where interference from neighboring
cells produces colored noise with a non-diagonal covariance
structure~\cite{zhao2025jsac}. Under colored noise, the standard
inversion-avoidance techniques used by AWGN-based methods lose their convergence
guarantees~\cite{wang2023bdnsepa}, and all known decentralized LMMSE solutions
for this regime require an explicit matrix inversion
 
The Block Coordinate Descent
MMSE (BCD-MMSE) algorithm was recently proposed~\cite{zhao2025jsac} as the first decentralized LMMSE equalizer to operate
under colored noise. In BCD-MMSE, each DU performs a local $M_c \times M_c$
Hermitian matrix inversion once per coherence block, reusing the result across
all received symbols. Despite its algorithmic completeness, no hardware
implementation of this inversion has been reported: the work
of~\cite{zhao2025jsac} provides neither hardware resource use, nor power, nor latency
figures for any target platform. Field-Programmable Gate Arrays (FPGAs) are
a promising target for this operation, offering higher energy efficiency and
lower latency than GPU-based solutions at the power budgets required by
distributed base-station nodes~\cite{jeon2019,li2019cd}.

This work proposes and evaluates an FPGA accelerator for the $M_c \times M_c$
Hermitian matrix inversion required by each DU in the BCD-MMSE algorithm. The
accelerator is based on Cholesky decomposition, an exact inversion method
applicable under a non-diagonal noise covariance. The contributions are
evaluated in terms of detection quality (Bit Error Rate - BER), resource usage
(LUT count and DSP blocks), and power consumption.

